% Auto-generated master file for t-test tables
% This file includes all benchmark t-test tables with configurable captions
%
% Usage: Include this file in your LaTeX document
% To customize captions, redefine the caption commands before including:
%   \renewcommand{\ttestCaptionBenchmarkName}{Your Custom Caption}
% To set the path to the table files (if not in same directory):
%   \renewcommand{\ttestBasePath}{path/to/tables/}
%
% Requires: booktabs, float, graphicx, currfile packages
%   \usepackage{booktabs}
%   \usepackage{float}
%   \usepackage{graphicx}
%   \usepackage{currfile}

% Base path for table files (empty = same directory)
\providecommand{\ttestBasePath}{\currfiledir}

\subsection{T-test Results for scenario 2: Parallel test execution performance}

% Legend explaining column abbreviations and environment names
{\footnotesize
\noindent\textbf{Metrics:} w = wall\_time\_seconds, c = system\_cpu\_delta\_from\_baseline, m = system\_memory\_delta\_mb\_max, hm = host\_peak\_vm\_memory\_mb.
Bold values indicate statistical significance (p < 0.05).

\noindent\textbf{Environments:} Windows 11 native
d... = Windows 11 native
default; Windows 11 native
p... = Windows 11 native
parallel (20) unittests; Debian 12 native
de... = Debian 12 native
default; Debian 12 native
pa... = Debian 12 native
parallel (20) unittests.
}

% Default caption definitions
\providecommand{\ttestCaptionScenarioUnittestsParallelUnitTests}{T-test p-values for parallel unit\_tests}

% Tables

\begin{table}[H]
\centering
\input{\ttestBasePath scenario_2_unittests_parallel_ttest_unit_tests.tex}
\caption{\ttestCaptionScenarioUnittestsParallelUnitTests}
\label{tab:ttest_scenario_2_unittests_parallel_unit_tests}
\end{table}